
\chapter{Not a Limit, but an Atlas}
\label{ch:singular-limit}

\section{The folklore and its discontents}

It is a piece of standard folklore, repeated in the first pages of countless
textbooks, that classical mechanics is ``what remains'' of relativity when
$c\to\infty$ (equivalently, when the characteristic speeds of the problem
satisfy $\beta = v/c \to 0$) and of quantum mechanics when $\hbar\to 0$. The
picture suggested by this way of speaking is a comforting one: there is a
single family of theories, continuously parametrized by $1/c$ and $\hbar$, with
special relativity and quantum mechanics sitting at generic points of the
family and classical mechanics sitting at the corner where both parameters
vanish. Classical mechanics, on this picture, is not really an autonomous
theory at all; it is a limit point, an image, a shadow cast by the deeper
theories onto the hyperplane $1/c=\hbar=0$.

This is too generous to classical mechanics and too glib about what a ``limit''
is. Two observations, already familiar to anyone who has tried to take the
folklore literally, show that something more interesting is going on.

First, the mathematical structures involved are not nested in the way the
picture requires. The phase space of a classical system and the Hilbert space
of its quantum counterpart are not the same kind of object related by a small
deformation; one is a symplectic manifold, the other a vector space with a
Hermitian form, and the passage between them (geometric quantization, or its
inverse) is a substantial piece of mathematics, not a Taylor expansion.
Likewise, the Galilei group is not a subgroup of the Poincar\'e group that one
recovers by ``turning off'' some generators; as we will see in
\S\ref{sec:contraction}, it is not even isomorphic, as an abstract structure,
to any subgroup or quotient of the Poincar\'e group. The relation between the
two is a \emph{contraction}, a specific and rather violent operation on Lie
algebras, not an inclusion.

Second, and more strikingly, classical mechanics contains objects that have no
counterpart at all on the other side of the limit, no matter how one takes it.
The rigid body is the standard example: an object whose configuration space is
finite-dimensional ($SE(3)$, or a quotient of it) precisely because every
internal distance is exactly conserved. Special relativity excludes such an
object outright, for any finite $c$, on causal grounds (\S\ref{sec:rigid}).
Friction and other dissipative forces are a second example: they break
time-reversal invariance and shrink phase-space volume, while the unitary
evolution of an isolated quantum system does neither, at any order in $\hbar$
(\S\ref{sec:friction}). These are not objects that quantum mechanics or
relativity describe badly, or only approximately; they are objects that
literally cannot be the image of anything on the quantum or relativistic side,
under any map deserving the name ``limit.''

The two observations reinforce each other. If the limit were regular --- if
solutions, symmetry groups, and observables all deformed continuously and were
connected by convergent, invertible correspondences --- then every classical
object would have to arise as the limit of \emph{something}, however
degenerate, on the quantum or relativistic side. The fact that some classical
objects have no such precursor is the symptom; the mismatch in mathematical
structure is the underlying disease. The right diagnosis, as is often noted, is
that these are \emph{singular} limits. But that word, by itself, is a label,
not an explanation --- and one of the aims of this chapter is to insist on the
difference, since a name that is not cashed out mechanistically explains
nothing more than does the word ``emergence.'' The other, larger aim is to
propose that once the mechanism behind ``singular'' is made precise in a few
cases, a different picture of classical mechanics suggests itself: not a limit
point at all, but something closer to an \emph{atlas} of local charts, each
valid in its own regime, glued together where their domains overlap, with no
single chart covering the whole territory.

\section{Two objects without a counterpart}

\subsection{Rigid bodies and the speed of sound}
\label{sec:rigid}

A perfectly rigid body is one all of whose interparticle distances are exactly
constant in time, in every frame simultaneously. Equivalently: a disturbance
applied at one point of the body is felt everywhere else instantaneously. This
is not merely difficult to arrange relativistically; it is forbidden by the
causal structure of Minkowski space at \emph{any} finite $c$, since it requires
a signal speed of $\infty$ inside the body. A real elastic body obeys, in the
simplest (longitudinal, one-dimensional) approximation, a wave equation
\begin{equation}
\pdv{^2 u}{t^2} = c_s^2\, \pdv{^2 u}{x^2}, \qquad c_s = \sqrt{E/\rho},
\end{equation}
with $E$ the elastic modulus and $\rho$ the density; the Newtonian rigid body
is recovered only in the further, separate idealization $c_s\to\infty$ of
\emph{this} equation, which turns a hyperbolic equation with finite propagation
speed into an elliptic constraint enforcing instantaneous equilibration.
Relativistic causality demands $c_s \le c$, so this second idealization is
unavailable once the first ($c$ finite) is taken seriously: a body ``as rigid
as possible'' in special relativity is precisely a body whose sound speed
equals $c$, not a body with no internal dynamics at all.

Max Born's 1909 attempt to define a relativistic notion of rigidity, and its
subsequent elaboration by Herglotz and Noether, confirms this from a different
direction.\footnote{M. Born, ``Die Theorie des starren K\"orpers in der
Kinematik des Relativit\"atsprinzips,'' \emph{Annalen der Physik}
\textbf{335}(11), 1--56 (1909); G. Herglotz, ``\"Uber den vom Standpunkt des
Relativit\"atsprinzips aus als starr zu bezeichnenden K\"orper,'' \emph{Annalen
der Physik} \textbf{336} (1910); F. Noether, \emph{Annalen der Physik}
\textbf{336} (1910).} Born's condition --- that each infinitesimal volume
element be undeformed in its instantaneous rest frame --- is the closest
relativistic analogue of rigidity one can write down without demanding infinite
signal speed. Even so, the Herglotz--Noether theorem shows that a body
satisfying it has only three kinematic degrees of freedom, not the six (three
translational, three rotational) of a Newtonian rigid body, and in particular
cannot be spun up from rest while remaining Born-rigid throughout --- this is
the content of the Ehrenfest paradox for the relativistic rotating disk. There
is, in other words, no relativistic object that specializes to the Newtonian
rigid body as $c\to\infty$; there is only a strictly smaller, strictly more
restrictive relativistic notion that happens to share Born's name with it.

Note the structure of the argument: it involves \emph{two} infinite-speed
idealizations that happen to coincide only in the fully non-relativistic,
classical regime --- the speed of light and the speed of sound in the medium.
Sending $c\to\infty$ alone, holding $c_s$ finite, gives ordinary
(non-relativistic) elasticity, not rigid-body mechanics. One needs a second,
independent limit, $c_s\to\infty$, to reach the rigid body, and that second
limit is available at all only once causality (bounded by $c$) has already been
relaxed. Rigid-body mechanics is thus not sitting ``one step away'' from
relativity along a single road; it sits at the intersection of two different
singular limits taken along two different roads.

\subsection{Friction and the closure of unitary evolution}
\label{sec:friction}

The evolution of an isolated quantum system is unitary, 
$\lvert\psi(t)\rangle =e^{-iHt/\hbar}\lvert\psi(0)\rangle$, 
for a time-independent Hamiltonian $H$.
This evolution is invertible, conserves the norm of the state, and --- through
the Ehrenfest and semiclassical limits discussed in \S\ref{sec:wkb} --- reduces
at leading order in $\hbar$ to ordinary Hamiltonian mechanics, which is itself
invertible and preserves phase-space volume (Liouville's theorem). Friction
does none of these things: it is dissipative, it privileges a direction of
time, and it shrinks the volume occupied by a bundle of trajectories in phase
space. No amount of expanding $e^{-iHt/\hbar}$ in powers of $\hbar$, for a
closed system with a time-independent, self-adjoint $H$, will ever produce a
term that does these things, because unitarity is exact at every order.

Friction is not, therefore, a phenomenon that the $\hbar\to 0$ limit of an
isolated quantum system approximates poorly; it is a phenomenon that lies
outside the entire category of closed, unitary evolutions, at any value of
$\hbar$. To obtain it, one must change the object under discussion, not just
take a limit of it: enlarge the system to include an environment, and discard
(trace out) the environment's degrees of freedom. The Caldeira--Leggett model
is the standard worked example --- a system bilinearly coupled to a large bath
of harmonic oscillators, whose exact elimination yields, in both the quantum
and the classical treatment, a generalized Langevin equation
\begin{equation}
m\ddot q(t) + \int_0^t \gamma(t-t')\,\dot q(t')\,\dd t' + V'(q) = \xi(t),
\end{equation}
with the memory kernel $\gamma$ and the noise correlator $\langle
\xi(t)\xi(t')\rangle$ tied together by a fluctuation--dissipation relation
fixed by the bath's spectral density.\footnote{A. O. Caldeira and A. J.
Leggett, ``Path Integral Approach to Quantum Brownian Motion,'' \emph{Physica
A} \textbf{121}(3), 587--616 (1983).} Friction, in other words, is reached only
by a limit taken in a \emph{different direction} from $\hbar\to 0$: many bath
degrees of freedom, weak coupling, short bath correlation time. It belongs to a
chart of the ``space of mechanical theories'' that the closed-system
semiclassical chart simply does not cover, at any resolution.

\section{What ``singular'' actually means}
\label{sec:singular}

To call these limits ``singular'' is, so far, only a name for our discomfort.
It becomes an explanation once we say, case by case, \emph{what fails to
converge and why}. Two mechanisms are illustrated below; they turn out to be
different mechanisms, which is itself worth noting; a family of ``singular''
phenomena is not the same thing as a single one wearing different clothes.

\subsection{The semiclassical limit as a singular limit}
\label{sec:wkb}

Write the wavefunction as 
$\psi(x,t) = A(x,t)\,e^{iS(x,t)/\hbar}$ 
and substitute into the Schr\"odinger equation 
$i\hbar\,\partial_t\psi = -\tfrac{\hbar^2}{2m}\nabla^2\psi + V\psi$.
Collecting powers of $\hbar$ gives, at order $\hbar^0$,
\begin{equation}
\pdv{S}{t} + \frac{(\nabla S)^2}{2m} + V = 0,
\end{equation}
the Hamilton--Jacobi equation for the classical action, and at order $\hbar^1$
a continuity equation 
$\partial_t(A^2) + \nabla\!\cdot(A^2\nabla S/m) = 0$
transporting the probability density along classical trajectories. This is the
correct and celebrated part of the folklore: to leading order, quantum
mechanics does reduce smoothly to classical mechanics in its Hamilton--Jacobi
form.

The reduction is singular, however, precisely where classical trajectories
focus or cross: at caustics and classical turning points, $A$ formally
diverges, several branches of $S$ become relevant at once, and the interference
between them does not disappear as $\hbar\to 0$ --- it is squeezed into a
shrinking neighbourhood of the caustic, but is never literally absent for any
$\hbar>0$.\footnote{See M. V. Berry, ``Singular Limits,'' \emph{Physics Today}
\textbf{55}(5), 10--11 (2002), whose simulations show the semiclassical
probability density varying smoothly except at isolated points where the phase
is undefined; and M. V. Berry, ``Asymptotics, Singularities and the Reduction
of Theories,'' in \emph{Logic, Methodology and Philosophy of Science IX}, eds.
D. Prawitz, B. Skyrms, D. Westerst\aa hl (Elsevier, 1994), 597--607.} Bridging
the oscillatory (classically allowed) and exponential (classically forbidden)
WKB solutions across a turning point requires a distinct local approximation
--- an Airy function --- together with connection formulas relating its two
asymptotic regimes back to the outer WKB solutions on either side. Those
connection formulas are not a decoration on the semiclassical limit; they are
the actual mathematical content of the reduction near its singular points, and
they are, formally, a transition map between two local charts.

\subsection{The non-relativistic limit as a group contraction}
\label{sec:contraction}

The Poincar\'e Lie algebra, with rotation generators $J_i$, boosts $K_i$,
momenta $P_i$, and the Hamiltonian $H$, closes as (up to conventional factors
of $i$)
\begin{align}
[J_i,J_j] &= i\hbar\,\epsilon_{ijk}J_k, & [J_i,K_j] &= i\hbar\,\epsilon_{ijk}K_k, & [K_i,K_j] &= -\frac{i\hbar}{c^2}\epsilon_{ijk}J_k,\\
[J_i,P_j] &= i\hbar\,\epsilon_{ijk}P_k, & [K_i,P_j] &= \frac{i\hbar}{c^2}\delta_{ij}H, & [K_i,H] &= i\hbar\,P_i,
\end{align}
with $[P_i,P_j]=[P_i,H]=0$. Sending $c\to\infty$ term by term kills the
right-hand sides of $[K_i,K_j]$ and $[K_i,P_j]$, so that boosts become abelian
and commute with momenta: this is exactly the kinematic content of the Galilei
algebra. \.{I}n\"on\"u and Wigner formalized this operation in 1953 under the
name \emph{group contraction} --- a singular change of basis, parametrized by
$1/c$, under which the structure constants of the algebra jump discontinuously
at $1/c=0$ rather than deforming smoothly through it.\footnote{E. \.{I}n\"on\"u
and E. P. Wigner, ``On the Contraction of Groups and Their Representations,''
\emph{Proc. Natl. Acad. Sci. U.S.A.} \textbf{39}(6), 510--524 (1953).}

But the naive contraction reproduces only part of what non-relativistic quantum
mechanics actually needs. The physically correct Galilei algebra has $[K_i,P_j]
= i\hbar\,\delta_{ij}\,m$, with the particle's mass $m$ appearing as a genuine
\emph{central charge} --- a new term with no counterpart anywhere in the
Poincar\'e algebra, since letting $H/c^2 \to \text{const.}$ as $c\to\infty$
while $H$ itself stays finite is not something the naive limit does on its own.
The mass term has to be inserted by hand, by first extending the Poincar\'e
algebra trivially with a central element (the Poincar\'e algebra itself admits
no non-trivial central extension) and only then contracting.\footnote{See,
e.g., the explicit construction in the ``generalized \.{I}n\"on\"u--Wigner
contraction'' of central-extended Poincar\'e to central-extended Galilei
algebras. The necessity of the central extension for the Galilei group's
\emph{projective} representations to be realized as genuine unitary
representations is Bargmann's theorem: V. Bargmann, ``On Unitary Ray
Representations of Continuous Groups,'' \emph{Annals of Mathematics}
\textbf{59}(1), 1--46 (1954).} Physically, this central charge is responsible
for the Galilean mass superselection rule --- no coherent superposition of
states of different mass is physically realizable non-relativistically --- a
structural feature of the non-relativistic theory that has no analogue on the
relativistic side, precisely because it is not present anywhere for $c<\infty$
and appears only in the limit itself, discontinuously, as new cohomology
available to the contracted algebra but not to the original one.

The moral of both examples is the same, though the mechanism differs. In
\S\ref{sec:wkb}, a smooth leading-order reduction coexists with genuinely new,
non-perturbative structure (caustics) concentrated at special points; in this
section, the leading-order reduction is not even smooth --- the limit has to be
patched, after the fact, with structure (the central extension) that the parent
algebra could not supply. Calling both of these ``singular limits'' is
convenient shorthand, not an explanation; the explanation is the mechanism, and
the mechanisms are different mechanisms.

\section{From a point to an atlas}

\subsection{Matched asymptotics are already an atlas}

The technology used to handle both examples above already has the structure
this chapter is arguing for. The oscillatory WKB solution valid away from a
turning point, and the Airy-function solution valid near it, are two local
approximations, each adequate on an open region of the relevant parameter space
and inadequate outside it; the connection formulas that relate them on their
overlap are exactly transition functions in the sense of differential geometry.
Boundary-layer methods in fluid mechanics --- an ``outer'' inviscid solution
away from a wall, an ``inner'' solution resolving the boundary layer, matched
in an intermediate region --- have precisely the same shape, and are the
paradigm case discussed under the heading of singular perturbation theory. This
is not a metaphor grafted onto existing mathematics after the fact; matched
asymptotic expansions simply \emph{are}, formally, an atlas of local charts
glued by transition maps, for the elementary reason that no single formula is
uniformly valid over the whole domain.

\subsection{Classical mechanics as an overlap, not an image}

The proposal, then, is to stop picturing ``the classical limit'' as a single
point approached along a single road in a single parent theory's parameter
space, and to picture it instead as a boundary region of a larger space of
mechanical theories, coordinatized by several independent small (or large)
parameters at once: the action ratio $\hbar/S$; the kinematic ratio $v/c$; the
ratio of an internal signal speed to $c$ (\S\ref{sec:rigid}); and a
coupling-to-environment parameter separating closed from open dynamics
(\S\ref{sec:friction}). Each of these, taken to its extreme, defines its own
local chart, valid within its own domain and glued to its neighbours by its own
matching conditions --- WKB connection formulas for $\hbar/S\to 0$, the
contraction of \S\ref{sec:contraction} for $v/c\to 0$, elastic-to-rigid
matching for the internal wave speed, Markovian reduction for the environment
coupling.

No single one of these charts contains everything we call classical mechanics.
The $\hbar\to 0$ chart alone gives Hamiltonian, non-dissipative, non-rigid
mechanics; it has to be supplemented by the $c_s\to\infty$ chart to admit rigid
bodies at all, and by the open-system chart to admit friction. ``Classical
mechanics,'' on this picture, is not the image of one map from one deeper
theory; it is the region covered by the union of several such charts, each
singular on its own boundary, overlapping in the regime we actually inhabit ---
moderate sizes, moderate speeds, weakly and unavoidably coupled to an
environment. And an object built by extending the mathematics of a single chart
past the region where it was ever matched to its neighbours --- an exactly
rigid lever transmitting force instantaneously; a closed, reversible
``explanation'' of a manifestly dissipative process --- is exactly the sort of
object this book has elsewhere called a paradox of mixing the real world with
the ideal one. It is a piece of one chart's bookkeeping mistaken for a feature
of the territory.

\section{Back to the real world}

This is, in the end, a modest correction to the folklore rather than a
rejection of it. Each of the limits discussed above is real, and each does
genuine reductive work within its own domain; the error lies only in imagining
that a single one of them, extended far enough, ought to hand back the whole of
classical mechanics, rigid bodies, friction, and all. It should not, any more
than a single map projection of the earth should be expected to preserve area,
angle, and distance all at once. What classical mechanics actually is, on this
reading, is the overlap of several distinct singular reductions, each exact in
its own chart and each silent about what happens in the others' --- which is
also, not coincidentally, why the objects that live only in the overlap of
several charts at once (a rigid body set gently down with friction on a table)
feel so unproblematic in practice and turn out, on inspection, to have no
single fundamental theory behind them at all. It is this same conflation of a
local, ideal-world chart with the whole of the real-world territory that we
take up object by object in the chapter that follows, beginning with Hertz's
forceless mechanics.

